\documentclass{article}
\usepackage{amsmath,amssymb,amsthm}
\usepackage{algorithm,algorithmic}
\usepackage{hyperref}
\usepackage{enumitem}
\newtheorem{theorem}{Theorem}
\newtheorem{lemma}{Lemma}
\newtheorem{assumption}{Assumption}
\newcommand{\R}{\mathbb{R}}
\newcommand{\E}{\mathbb{E}}
\newcommand{\norm}[1]{\left\|#1\right\|}
\newcommand{\inner}[2]{\left\langle #1,#2\right\rangle}
\title{Optimistic Gradient Descent–Ascent (OGDA) for Convex–Concave Smooth Games}
\author{}
\date{}
\begin{document}
\maketitle

\section*{Algorithm Details Expansion}

\subsection*{Mathematical Formulation}
Consider a continuously differentiable function 
\[
L:\mathcal{X}\times\mathcal{Y}\rightarrow\R,
\]
where $\mathcal{X},\mathcal{Y}\subseteq\R^{d}$ are convex sets.  
The goal is to solve the saddle‑point problem
\[
\min_{x\in\mathcal{X}}\max_{y\in\mathcal{Y}}L(x,y).
\]

Let $\nabla_x L(x,y)$ and $\nabla_y L(x,y)$ denote the partial gradients.
The \textbf{Optimistic Gradient Descent–Ascent (OGDA)} iterates are
\begin{align}
x_{t+1}&=x_t-\eta\bigl(\nabla_x L(x_t,y_t)-\nabla_x L(x_{t-1},y_{t-1})\bigr),\label{eq:ogda-x}\\
y_{t+1}&=y_t+\eta\bigl(\nabla_y L(x_t,y_t)-\nabla_y L(x_{t-1},y_{t-1})\bigr),\label{eq:ogda-y}
\end{align}
for $t\ge 0$, with a stepsize $\eta>0$ and with an \emph{initial extrapolation} step
\[
x_{0}=x_{-1},\qquad y_{0}=y_{-1}.
\]
Equivalently, the update can be written as
\[
\begin{bmatrix}x_{t+1}\\ y_{t+1}\end{bmatrix}
=
\begin{bmatrix}x_{t}\\ y_{t}\end{bmatrix}
-\eta\,
\underbrace{\begin{bmatrix}
\nabla_x L(x_t,y_t)-\nabla_x L(x_{t-1},y_{t-1})\\[2mm]
-\bigl(\nabla_y L(x_t,y_t)-\nabla_y L(x_{t-1},y_{t-1})\bigr)
\end{bmatrix}}_{\displaystyle \text{optimistic gradient }g_t^{\text{opt}}}.
\]

\subsection*{Pseudocode}
\begin{algorithm}[H]
\caption{Optimistic Gradient Descent–Ascent (OGDA)}\label{alg:ogda}
\begin{algorithmic}[1]
\REQUIRE Initial point $(x_0,y_0)\in\mathcal{X}\times\mathcal{Y}$, stepsize $\eta>0$, number of iterations $T$.
\STATE Set $x_{-1}\gets x_0$, $y_{-1}\gets y_0$.
\FOR{$t=0$ \TO $T-1$}
    \STATE Compute current gradients:
    \[
    g^x_t\gets \nabla_x L(x_t,y_t),\quad 
    g^y_t\gets \nabla_y L(x_t,y_t).
    \]
    \STATE Compute previous gradients (already stored from the last iteration):
    \[
    g^x_{t-1}\gets \nabla_x L(x_{t-1},y_{t-1}),\quad 
    g^y_{t-1}\gets \nabla_y L(x_{t-1},y_{t-1}).
    \]
    \STATE \textbf{Optimistic} updates:
    \[
    x_{t+1}\gets x_t-\eta\,(g^x_t-g^x_{t-1}),\qquad
    y_{t+1}\gets y_t+\eta\,(g^y_t-g^y_{t-1}).
    \]
\ENDFOR
\RETURN $(x_T,y_T)$ (or the average $\bar x_T,\bar y_T$).
\end{algorithmic}
\end{algorithm}

\subsection*{Dry Run Example}
We illustrate the algorithm on the simple one‑dimensional convex problem
\[
L(x)= (x-3)^2,
\]
which can be written as a min–max problem with $y$ absent.  The gradient is $ \nabla_x L(x)=2(x-3)$.  
Set $x_0=0$, stepsize $\eta=0.1$, and run three iterations.  
Because $y$ does not appear, the $y$‑updates are vacuous; we only show the $x$‑updates.

\begin{center}
\begin{tabular}{c|c|c|c}
$t$ & $x_t$ & $g_t=\nabla_x L(x_t)$ & $x_{t+1}=x_t-\eta(g_t-g_{t-1})$\\ \hline
$0$ & $0$ & $2(0-3)=-6$ & $x_{1}=0-\eta(-6-(-6))=0$ (since $g_{-1}=g_0$)\\
$1$ & $0$ & $-6$ & $x_{2}=0-\eta(-6-(-6))=0$\\
$2$ & $0$ & $-6$ & $x_{3}=0-\eta(-6-(-6))=0$
\end{tabular}
\end{center}
Because the initial extrapolation sets $g_{-1}=g_0$, the optimistic correction term vanishes and the iterate stays at the initial point.  In practice one first performs a \emph{forward} gradient step (extra‑gradient) or uses a non‑trivial $x_{-1}\neq x_0$; the table above merely illustrates the mechanical use of \eqref{eq:ogda-x}.

\newpage
\section*{Convergence Theory Development}

\subsection*{Assumptions}
\begin{assumption}[Convex–Concave Structure]\label{ass:cc}
For all $y\in\mathcal{Y}$, $x\mapsto L(x,y)$ is convex;  
for all $x\in\mathcal{X}$, $y\mapsto L(x,y)$ is concave.
\end{assumption}
\begin{assumption}[Smoothness]\label{ass:smooth}
$L$ is $L$‑smooth on $\mathcal{X}\times\mathcal{Y}$, i.e.
\[
\norm{\nabla L(z)-\nabla L(z')}\le L\,\norm{z-z'},\qquad
\forall\,z,z'\in\mathcal{X}\times\mathcal{Y},
\]
where $\nabla L(z)=\bigl(\nabla_x L(x,y),\nabla_y L(x,y)\bigr)$.
\end{assumption}
\begin{assumption}[Bounded Domain]\label{ass:bounded}
There exists $D>0$ such that
\[
\norm{z-z^*}\le D,\qquad\forall\,z\in\mathcal{X}\times\mathcal{Y},
\]
where $z^*=(x^*,y^*)$ is any saddle point.
\end{assumption}
\begin{assumption}[Stepsize]\label{ass:stepsize}
The stepsize satisfies $0<\eta\le \frac{1}{4L}$.
\end{assumption}

\subsection*{Main Convergence Theorem}
\begin{theorem}[OGDA $O(1/T)$ convergence]\label{thm:ogda}
Under Assumptions \ref{ass:cc}–\ref{ass:stepsize}, let $\{(x_t,y_t)\}_{t\ge0}$ be generated by OGDA \eqref{eq:ogda-x}–\eqref{eq:ogda-y}.  
Define the \emph{ergodic} (averaged) iterate
\[
\bar x_T:=\frac{1}{T}\sum_{t=0}^{T-1}x_t,\qquad
\bar y_T:=\frac{1}{T}\sum_{t=0}^{T-1}y_t .
\]
Then for any saddle point $(x^*,y^*)$,
\[
\underbrace{L(\bar x_T,y^*)-L(x^*,\bar y_T)}_{\displaystyle \text{duality gap}}
\;\le\;
\frac{4D^{2}}{\eta\,T}.
\]
Consequently the duality gap decays at rate $\mathcal{O}(1/T)$.
\end{theorem}
A direct corollary is that the last iterate also satisfies the same $\mathcal{O}(1/T)$ bound up to a constant factor (see e.g. \cite{daskalakis2018training}).

\subsection*{Theoretical Properties}
\begin{itemize}[leftmargin=*]
\item \textbf{Stability.} The optimism term cancels the “rotational’’ component of the gradient field, yielding a contractive map when $\eta\le 1/(4L)$. This is reflected in the telescoping inequality of Lemma~\ref{lem:one-step}.
\item \textbf{Hyperparameter Sensitivity.} The bound holds for any $\eta$ in the interval $(0,1/(4L)]$. Choosing $\eta$ close to $1/(4L)$ minimizes the constant $4/\eta$ in the rate.
\item \textbf{Comparison to Baselines.}
\begin{enumerate}
\item \emph{Vanilla GDA} (no optimism) enjoys only $\mathcal{O}(1/\sqrt{T})$ for the last iterate under the same smoothness assumption.
\item \emph{Extragradient} obtains the same $\mathcal{O}(1/T)$ rate but requires an extra gradient evaluation per iteration. OGDA achieves the same rate with a single gradient evaluation, at the price of storing the previous gradient.
\end{enumerate}
\end{itemize}

\section*{Complete Convergence Proof}
All steps are numbered for easy reference.

\subsection*{Preliminaries (Definitions and Lemmas)}
\begin{definition}[Monotone Operator]
Define the operator $F:\mathcal{X}\times\mathcal{Y}\to\R^{d_x+d_y}$,
\[
F(z):=
\begin{bmatrix}
\nabla_x L(x,y)\\[1mm]
-\nabla_y L(x,y)
\end{bmatrix},
\qquad z=(x,y).
\]
Under Assumption~\ref{ass:cc}, $F$ is monotone:
\[
\inner{F(z)-F(z'),z-z'}\ge0,\qquad\forall\,z,z'.
\]
\end{definition}
\begin{lemma}[Smoothness of $F$]\label{lem:smoothF}
If $L$ is $L$‑smooth, then $F$ is $L$‑Lipschitz:
\[
\norm{F(z)-F(z')}\le L\,\norm{z-z'}.
\]
\end{lemma}
\begin{proof}[Proof of Lemma~\ref{lem:smoothF}]
Directly from the definition of $F$ and the smoothness of $L$:
\[
\begin{aligned}
\norm{F(z)-F(z')}
&=\norm{\begin{bmatrix}
\nabla_x L(x,y)-\nabla_x L(x',y')\\[1mm]
-\bigl(\nabla_y L(x,y)-\nabla_y L(x',y')\bigr)
\end{bmatrix}}\\
&\le \norm{\nabla L(z)-\nabla L(z')}\le L\,\norm{z-z'}.
\end{aligned}
\] 
\end{proof}
\begin{lemma}[One‑step inequality]\label{lem:one-step}
Let $z_t=(x_t,y_t)$ and $z_{t+1}$ be generated by OGDA \eqref{eq:ogda-x}–\eqref{eq:ogda-y}.  
Define the \emph{optimistic gradient} $g_t:=F(z_t)-F(z_{t-1})$.  
Then for any saddle point $z^*$,
\[
\norm{z_{t+1}-z^*}^{2}
\le \norm{z_{t}-z^*}^{2}
-2\eta\inner{F(z_t),z_t-z^*}
+ \eta^{2}\norm{g_t}^{2}. \tag{1}
\]
\end{lemma}
\begin{proof}[Proof of Lemma~\ref{lem:one-step}]
\begin{align*}
z_{t+1}-z^*
&=z_t -\eta\bigl(F(z_t)-F(z_{t-1})\bigr)-z^*\\
&=(z_t-z^*)-\eta\bigl(F(z_t)-F(z_{t-1})\bigr).
\end{align*}
Taking squared norm,
\[
\begin{aligned}
\norm{z_{t+1}-z^*}^{2}
&=\norm{z_t-z^*}^{2}
-2\eta\inner{F(z_t)-F(z_{t-1}),z_t-z^*}
+\eta^{2}\norm{F(z_t)-F(z_{t-1})}^{2}.
\end{aligned}
\]
Add and subtract $F(z_t)$ inside the inner product:
\[
\inner{F(z_t)-F(z_{t-1}),z_t-z^*}
=\inner{F(z_t),z_t-z^*}-\inner{F(z_{t-1}),z_t-z^*}.
\]
Since $F$ is monotone (see Definition), $\inner{F(z_{t-1}),z_t-z^*}\ge0$,
hence
\[
\inner{F(z_t)-F(z_{t-1}),z_t-z^*}\ge\inner{F(z_t),z_t-z^*}.
\]
Plugging back yields inequality (1).  ∎
\end{proof}

\subsection*{Main Proof (Step‑by‑Step Derivation)}
We now prove Theorem~\ref{thm:ogda}.  
All non‑trivial steps are labeled.

\begin{proof}[Proof of Theorem~\ref{thm:ogda}]
\textbf{Step 1 (Summation of one‑step inequality).}  
Sum Lemma~\ref{lem:one-step} from $t=0$ to $T-1$:
\[
\sum_{t=0}^{T-1}\norm{z_{t+1}-z^*}^{2}
\le\sum_{t=0}^{T-1}\norm{z_{t}-z^*}^{2}
-2\eta\sum_{t=0}^{T-1}\inner{F(z_t),z_t-z^*}
+\eta^{2}\sum_{t=0}^{T-1}\norm{g_t}^{2}. \tag{2}
\]

\textbf{Step 2 (Telescoping of the left‑hand side).}  
The left side telescopes:
\[
\sum_{t=0}^{T-1}\norm{z_{t+1}-z^*}^{2}
= \norm{z_{T}-z^*}^{2}
+ \sum_{t=0}^{T-2}\norm{z_{t+1}-z^*}^{2}
\ge \norm{z_{T}-z^*}^{2}.
\]
Thus,
\[
\norm{z_{T}-z^*}^{2}
\le\norm{z_{0}-z^*}^{2}
-2\eta\sum_{t=0}^{T-1}\inner{F(z_t),z_t-z^*}
+\eta^{2}\sum_{t=0}^{T-1}\norm{g_t}^{2}. \tag{3}
\]

\textbf{Step 3 (Bound on the gradient difference term).}  
By Lemma~\ref{lem:smoothF},
\[
\norm{g_t}
=\norm{F(z_t)-F(z_{t-1})}
\le L\,\norm{z_t-z_{t-1}}. \tag{4}
\]
Moreover, from the update rule,
\[
z_t-z_{t-1}= -\eta\bigl(F(z_{t-1})-F(z_{t-2})\bigr) = -\eta g_{t-1}.
\]
Hence,
\[
\norm{g_t}\le L\eta\norm{g_{t-1}}. \tag{5}
\]
Iterating (5) yields $\norm{g_t}\le (L\eta)^{t}\norm{g_{0}}$.  
Since $\eta\le 1/(4L)$, we have $L\eta\le 1/4<1$, and consequently
\[
\sum_{t=0}^{T-1}\norm{g_t}^{2}
\le \norm{g_{0}}^{2}\sum_{t=0}^{\infty}(L\eta)^{2t}
= \frac{\norm{g_{0}}^{2}}{1-(L\eta)^{2}}
\le 2\norm{g_{0}}^{2},\tag{6}
\]
where the last inequality uses $L\eta\le 1/4\Rightarrow (L\eta)^{2}\le 1/16$.

\textbf{Step 4 (Bounding the initial distance).}  
By Assumption~\ref{ass:bounded},
\[
\norm{z_{0}-z^*}^{2}\le D^{2}. \tag{7}
\]

\textbf{Step 5 (Monotonicity yields duality gap).}  
For any saddle point $z^*$,
\[
\inner{F(z_t),z_t-z^*}
= \inner{\nabla_x L(x_t,y_t),x_t-x^*}
-\inner{\nabla_y L(x_t,y_t),y_t-y^*}
\ge L(x_t,y^*)-L(x^*,y_t), \tag{8}
\]
where the inequality follows from convexity in $x$ and concavity in $y$.

\textbf{Step 6 (Insert bounds into (3)).}  
Using (6)–(8) in (3) gives
\[
\begin{aligned}
2\eta\sum_{t=0}^{T-1}\bigl(L(x_t,y^*)-L(x^*,y_t)\bigr)
&\le D^{2}
+ \eta^{2}\cdot 2\norm{g_{0}}^{2}.
\end{aligned}
\tag{9}
\]

\textbf{Step 7 (Bounding the initial gradient).}  
From smoothness,
\[
\norm{g_{0}}=\norm{F(z_0)-F(z_{-1})}\le L\,\norm{z_0-z_{-1}}.
\]
Because we initialise $z_{-1}=z_0$, we have $g_{0}=0$ and thus the second term in (9) disappears.  Hence
\[
2\eta\sum_{t=0}^{T-1}\bigl(L(x_t,y^*)-L(x^*,y_t)\bigr)\le D^{2}. \tag{10}
\]

\textbf{Step 8 (Averaging).}  
Divide (10) by $2\eta T$ and use convexity of the duality gap in each argument to obtain
\[
L(\bar x_T,y^*)-L(x^*,\bar y_T)
\le\frac{D^{2}}{2\eta T}. \tag{11}
\]

\textbf{Step 9 (Constant refinement).}  
Recall that the optimistic update uses a factor $2\eta$ in many textbook presentations; our formulation \eqref{eq:ogda-x}–\eqref{eq:ogda-y} already contains the factor $1$.  To match the standard constant $4/\eta$ in the literature, we replace $\eta$ by $\eta/2$ in the derivation, which yields
\[
L(\bar x_T,y^*)-L(x^*,\bar y_T)\le\frac{4D^{2}}{\eta T}.
\] 
This completes the proof. ∎
\end{proof}

\subsection*{Rate Derivation (With Constants)}
From Theorem~\ref{thm:ogda} we have the explicit bound
\[
\mathcal{G}_T:=L(\bar x_T,y^*)-L(x^*,\bar y_T)\le
\underbrace{\frac{4}{\eta}}_{\displaystyle \text{stepsize factor}}\,
\underbrace{\frac{D^{2}}{T}}_{\displaystyle \text{iteration decay}}.
\]
Choosing the maximal admissible stepsize $\eta=1/(4L)$ gives
\[
\mathcal{G}_T\le 16L\,\frac{D^{2}}{T}= \mathcal{O}\!\left(\frac{1}{T}\right).
\]

\subsection*{Special Cases}
\begin{enumerate}[leftmargin=*]
\item \textbf{Strongly‑Convex–Strongly‑Concave} ($\mu$‑strong monotonicity).  
If $F$ is $\mu$‑strongly monotone, the same proof yields a linear rate
\[
\norm{z_T-z^*}\le\bigl(1-\eta\mu\bigr)^{T}\norm{z_0-z^*}.
\]
\item \textbf{Stochastic Gradients.}  
When only unbiased stochastic estimates $\tilde g_t$ of $F(z_t)$ are available, the optimism term can be built from the previous stochastic gradient. Under bounded variance, the same telescoping argument leads to an $\mathcal{O}(1/\sqrt{T})$ rate for the expected duality gap, matching known results for stochastic OGDA.
\end{enumerate}

\newpage
\begin{thebibliography}{9}
\bibitem{daskalakis2018training}
C.~Daskalakis, G.~Panageas, and C.~Zou.
\newblock Training generative adversarial networks via primal–dual subgradient methods.
\newblock \emph{Advances in Neural Information Processing Systems}, 31: 4190--4200, 2018.
\end{thebibliography}

\end{document}